% authority: science_draft
\documentclass[11pt]{article}

\usepackage[margin=1in]{geometry}
\usepackage{amsmath,amssymb,amsthm}
\usepackage{booktabs}
\usepackage{enumitem}
\usepackage{longtable}

\newcommand{\Bcur}{B_{\mathrm{current}}}
\newcommand{\MSBridge}{\mathrm{MSStableMatterSemanticsBridge}_{v1}}
\newcommand{\SMSAR}{\mathrm{SourceMatterSemanticsAdoptionReadinessLaw}_{v1}}
\newcommand{\PMSTarget}{\mathrm{PositiveSourceMatterSemanticsTarget}_{v1}}
\newcommand{\PMSProf}{\mathrm{PositiveMSProfile}_{v1}}
\newcommand{\AdmSrc}{\mathrm{Adm}^{\mathrm{MS}}_{\mathrm{src}}}
\newcommand{\AdmPMS}{\mathrm{Adm}^{\mathrm{PMS}}_{\mathrm{src}}}
\newcommand{\RRE}{\mathrm{RR}_E}
\newcommand{\RRTarget}{\mathrm{RR}_{E}\mathrm{TheoremTarget}_{v1}}
\newcommand{\Qsrc}{Q^{\RRE}_{\mathrm{src}}}
\newcommand{\EqRR}{\equiv^{\RRE}_{\mathrm{src}}}
\newcommand{\SepRR}{\#^{\RRE}_{\mathrm{src}}}
\newcommand{\IrrRR}{\mathrm{Irr}^{\RRE}_{\mathrm{src}}}
\newcommand{\TrRR}{\mathrm{Tr}^{\RRE}_{\mathrm{src}}}
\newcommand{\Bottom}{\bot}
\newcommand{\Obstruction}{\mathsf{obstruction}}
\newcommand{\RRUnder}{\mathsf{RR\_E\_underdetermination\_obstruction}}
\newcommand{\MetricData}{\mathrm{MetricData}}
\newcommand{\geff}{g_{\mathrm{eff}}}

\newtheorem{definition}{Definition}
\newtheorem{condition}{Condition}
\newtheorem{target}{Theorem target}
\newtheorem{example}{Finite/local example}
\newtheorem{failure}{Fail-closed branch}
\newtheorem{nonconclusion}{Non-conclusion}

\title{RR\_E Theorem Target Formalizer v1}
\author{Ontology Formalizer Packet RT-20260701-022}
\date{2026-07-01}

\begin{document}
\maketitle

\section{Control Status}

This artifact records the fused Ontology Formalizer output for
\texttt{RT-20260701-022}, implementing v13 P5-T01. It defines
\[
  \RRTarget
\]
as a \emph{proposal-only} and \emph{draft/control} source-side theorem target.
The artifact answers what an \(\RRE\) record is, what it means to preserve
\(\RRE\) separation, what source-side irrelevance would require, which
assumptions are allowed, which assumptions would smuggle detector semantics,
which finite/local examples separate preservation from irrelevance, and what
fail-closed obstruction applies if irrelevance is underdetermined.

This packet does not prove a general \(\RRE\) irrelevance theorem. It does not
edit canonical ontology TeX. It does not create or adopt a source law. It does
not adopt \(\PMSProf\). It does not adopt \(\SMSAR\) as law. It does not adopt
source-extension data beyond exact scoped Gate Chair results. It does not
adopt \(\MetricData(E)\), change \(\geff\), adopt a coupling law, adopt matter
semantics, adopt detector semantics, derive or adopt matter coupling, import
stress-energy semantics, construct a stress-energy tensor, import a matter
action, derive Einstein equations, promote benchmark status, issue benchmark
closure, claim completed derivation, claim future source-extension
impossibility, or reject the global theory.

\begin{verbatim}
candidate_status: "proposal-only"
packet_payload_mode: "rr_e_theorem_target"
RR_E_irrelevance_theorem_proved: false
source_law_adopted: false
PositiveMSProfile_v1_adopted: false
SourceMatterSemanticsAdoptionReadinessLaw_v1_adopted_as_law: false
source_extension_data_adopted_beyond_exact_scoped_gate_result: false
matter_semantics_adopted: false
detector_semantics_adopted: false
coupling_law_adopted: false
matter_coupling_derived: false
matter_coupling_adopted: false
stress_energy_semantics_imported: false
stress_energy_tensor_constructed: false
matter_action_imported: false
MetricData_E_adopted: false
g_eff_scope_changed: false
Einstein_equations_derived: false
benchmark_promoted: false
completed_derivation_claimed: false
\end{verbatim}

\section{Question}

The narrow P5-T01 question is:
\[
\text{What source-side theorem target decides whether distinct \(\RRE\)
records may be identified}
\]
\[
\text{without importing detector semantics or target physics?}
\]

Under current authority, the answer can only be a theorem-target grammar,
proof-obligation surface, finite/local example family, and fail-closed
obstruction label. The theorem attempt or obstruction belongs to P5-T02.

\section{Explicitly Excluded Imports}

The theorem target forbids the following as mathematical premises:
\begin{itemize}[leftmargin=*]
  \item detector protocols, measurement outcomes, empirical calibration, and
  detector distinguishability;
  \item stress-energy semantics, stress-energy tensor, matter action,
  variational principle, and conservation law by imported theory;
  \item target metric, Lorentzian signature, target topology, target atlas,
  target coordinate transition, proper time, and target-geometric
  equivalence;
  \item \(\MetricData(E)\), \(\geff\) scope expansion, Einstein equations,
  benchmark fit, benchmark recovery, and process authority;
  \item registry metadata, validator status, handoff status, approval state,
  role identity, commit state, generated derivatives, and local cache state as
  mathematical evidence.
\end{itemize}

\section{Source-Side RR\_E Records}

\begin{definition}[Source-side \(\RRE\) record]
Let \(r\in\AdmSrc(E)\) be an admissible finite/local source record already
available to the P4 source-profile lane. Its retained \(\RRE\) component is a
source syntax class
\[
  \Qsrc(r).
\]
\(\Qsrc(r)\) is a record of retained source readout-response syntax. It is not
a detector protocol, not a measurement event, not an empirical calibration,
not stress-energy data, not a matter action, not \(\MetricData(E)\), not
\(\geff\), and not benchmark evidence.
\end{definition}

\begin{definition}[RR\_E transport certificate]
A source transport certificate
\[
  \TrRR(r,r')
\]
is a named source-side datum proving that the retained syntax classes
\(\Qsrc(r)\) and \(\Qsrc(r')\) may be identified inside the declared source
formalism. It cannot be supplied by detector semantics, target geometry,
benchmark status, process authority, or generated derivatives.
\end{definition}

\begin{definition}[RR\_E separation]
For \(r,r'\in\AdmSrc(E)\), write \(r\SepRR r'\) when
\[
  \Qsrc(r)\ne \Qsrc(r')
\]
and no source transport certificate \(\TrRR(r,r')\) is available. Separation
is a source-side nonidentification condition. It is not detector
distinguishability and not empirical calibration.
\end{definition}

\begin{definition}[RR\_E source equivalence]
For \(r,r'\in\AdmSrc(E)\), write \(r\EqRR r'\) only when a source transport
certificate or a source-only equivalence proof identifies their retained
\(\RRE\) syntax classes. Any proposed equivalence that depends on a forbidden
import returns \(\Obstruction\).
\end{definition}

\section{Separation Preservation And Irrelevance}

\begin{target}[RR\_E separation preservation]
For admissible records \(r,r'\), if \(r\SepRR r'\), then any later
candidate, quotient, or theorem attempt that identifies \(r\) and \(r'\) must
exhibit a source-only \(\EqRR\) proof. Without that proof, identification
returns \(\RRUnder\).
\end{target}

\begin{target}[RR\_E source-side irrelevance]
Let \(F\) be a future source-side object under test, such as a profile label,
candidate relation, or obstruction predicate. Source-side \(\RRE\)
irrelevance for \(F\) means:
\[
  r \sim_{\mathrm{src}} r'
  \ \text{outside the retained \(\RRE\) component}
  \quad \Longrightarrow \quad
  F(r)=F(r')
\]
for every admitted finite/local \(\RRE\) variation or quotient, with the proof
using only source-side definitions, source certificates, and source
transports. The statement is a theorem target, not a theorem proved here.
\end{target}

\begin{condition}[Allowed assumptions]
P5-T02 may use the following assumptions when attempting the theorem or
obstruction:
\begin{enumerate}[label=A\arabic*.,leftmargin=*]
  \item finite/local source records \(r\in\AdmSrc(E)\) or
  \(r\in\AdmPMS(E)\);
  \item source support syntax, source probe-token syntax, source certificate
  data, retained \(\RRE\) syntax, and no-target certificates;
  \item explicit source transports and finite source variations;
  \item scoped \(\MSBridge\), \(\SMSAR\), and \(\PMSProf\) evidence only as
  source-extension evidence/precondition context, never as adoption;
  \item prior P4 stress result that \(\RRE\) collapse without source transport
  fails closed.
\end{enumerate}
\end{condition}

\begin{condition}[Forbidden assumptions]
P5-T02 may not use detector protocol, empirical calibration, target geometry,
target metric, stress-energy object, matter action, imported conservation law,
\(\MetricData(E)\), \(\geff\) expansion, Einstein equations, benchmark fit,
generated derivative, registry row, validator result, handoff wording, role
identity, approval state, commit state, file order, or local cache state as a
premise of \(\EqRR\), \(\SepRR\), or \(\IrrRR\).
\end{condition}

\section{Finite/Local Examples}

\begin{example}[Separation without irrelevance]
Let \(r_0,r_1\in\AdmPMS(E)\) agree on all source components except
\[
  \Qsrc(r_0)=Q_0,\qquad \Qsrc(r_1)=Q_1,\qquad Q_0\ne Q_1,
\]
and suppose no \(\TrRR(r_0,r_1)\) is supplied. Then \(r_0\SepRR r_1\). A
future profile, label, or route predicate must preserve the distinction or
return \(\RRUnder\). It may not identify \(r_0\) and \(r_1\) by saying that no
detector could distinguish them.
\end{example}

\begin{example}[Equivalence by source transport]
Let \(r_0,r_1\in\AdmPMS(E)\) again differ only in retained \(\RRE\) syntax,
but suppose a named source transport certificate \(\tau\) proves
\[
  \tau(Q_0)=Q_1
\]
inside the declared source formalism. Then P5-T02 may attempt an \(\EqRR\)
proof. The proof must remain source-only and must not appeal to detector
responses or target physics.
\end{example}

\begin{example}[Irrelevance target]
Let \(F\) be a future source-side label or obstruction predicate. The
\(\RRE\)-irrelevance target for \(F\) requires proving
\[
  F(r_0)=F(r_1)
\]
for every admitted pair \(r_0,r_1\) related by \(\RRE\)-variation while all
non-\(\RRE\) source components are held fixed or transported. Showing this for
a single pair is only an example; the theorem target requires a specified
domain and quantifier.
\end{example}

\begin{example}[Underdetermination obstruction]
If neither a source transport certificate nor a source-only invariance proof
is available, current ontology does not derive the identification rule needed
for \(\RRE\) irrelevance. The correct result is
\(\RRUnder\), not detector semantics and not an impossibility theorem.
\end{example}

\section{Fail-Closed Obstruction Family}

\begin{failure}[RR\_E collapse without source proof]
If a candidate identifies \(r,r'\) with \(r\SepRR r'\) and supplies no
source-only \(\EqRR\) proof, the branch returns \(\RRUnder\).
\end{failure}

\begin{failure}[Detector-semantic import]
If \(\EqRR\), \(\SepRR\), or \(\IrrRR\) depends on detector protocols,
measurement semantics, empirical calibration, or detector distinguishability,
the branch returns \(\Obstruction\).
\end{failure}

\begin{failure}[Target or downstream import]
If \(\EqRR\), \(\SepRR\), or \(\IrrRR\) depends on target metric, target
atlas, stress-energy semantics, stress-energy tensor, matter action,
\(\MetricData(E)\), \(\geff\), Einstein equations, benchmark fit, or completed
derivation status, the branch returns \(\Obstruction\).
\end{failure}

\begin{failure}[Process-authority import]
If registry metadata, validator status, role identity, handoff wording,
approval state, commit state, generated derivatives, file order, or local
cache state is used as a premise of \(\EqRR\), \(\SepRR\), or \(\IrrRR\), the
branch returns \(\Obstruction\).
\end{failure}

\section{Proof Obligations for P5-T02}

P5-T02 must do at least one of the following:
\begin{enumerate}[label=O\arabic*.,leftmargin=*]
  \item prove a source-side \(\RRE\) irrelevance theorem on a declared
  finite/local domain;
  \item prove a source-side separation theorem showing nonidentification under
  precise hypotheses;
  \item construct a finite countermodel or witness where separation matters;
  \item record a precise \(\RRUnder\) obstruction with the missing source-side
  rule named explicitly;
  \item recommend route freeze only if repeated-burden or scoped-obstruction
  criteria are actually met.
\end{enumerate}

\section{Distance-to-GR Status}

\begin{longtable}{p{0.30\linewidth}p{0.58\linewidth}}
\toprule
Burden & Status after this packet \\
\midrule
Source ontology primitives & Unchanged; no canonical ontology edit. \\
Source equivalence EqSrc & Unchanged; local \(\RRE\) target only. \\
RetainH & Unchanged. \\
GenH & Unchanged. \\
ObsLoc\_lc & Prior source-localization context only. \\
Resp\_lc & Prior accepted response/localization context only. \\
M\_src & Unchanged; scoped source-only context retained. \\
\(g_{\mathrm{eff}}\) & Unchanged; no scope expansion. \\
Matter coupling & \(\RRTarget\) formalized as proposal-only theorem target;
no matter coupling. \\
Einstein equations & Blocked; no field-equation derivation. \\
Finite-variation robustness & Future P5-T02 obligation only. \\
Benchmark promotion & Blocked; no benchmark promotion. \\
Gate Chair status & Prior scoped evidence/precondition results retained; no
new gate verdict. \\
Current route freeze or hard-fail status & Not frozen; P5-T02 theorem attempt
or obstruction selected. \\
\bottomrule
\end{longtable}

\section{Non-Conclusions}

\begin{nonconclusion}[No theorem proof]
This artifact does not prove \(\RRE\) irrelevance or separation as a final
theorem. It defines theorem targets and obstruction labels.
\end{nonconclusion}

\begin{nonconclusion}[No adoption or promotion]
This artifact does not adopt source law, \(\PMSProf\), \(\SMSAR\) as law,
matter semantics, detector semantics, coupling law, matter coupling,
stress-energy semantics, matter action, \(\MetricData(E)\), \(\geff\),
Einstein equations, benchmark status, or completed derivation.
\end{nonconclusion}

\section{APA 7 Source Notes}

The AEther-Flow Research Project. (2026, July 1). \textit{PositiveSourceMatterSemanticsTarget
v1} [Draft/control TeX artifact].

The AEther-Flow Research Project. (2026, July 1). \textit{Positive Source
Matter-Semantics Profile Candidate v1} [Draft/control TeX artifact].

The AEther-Flow Research Project. (2026, July 1). \textit{PositiveMSProfile v1
Refuter Stress} [Draft/control TeX artifact].

\end{document}
